\chapter{Ground}
% Covers: outline A1-A5: atoms, the fencepost, earned arithmetic, the one selection, naming
% Chapter language (per Ch1 protocol): grammar. Arrow: lambda-calculus -> BNF.

\section{The atoms}
% Node-map: A1 (atoms {0,1,2}; 3 = the fencepost, a Limit approached never landed;
% ladder 0=Fact/1=Number/2=Natural/3=Limit as interpretation riding on the construction).
% Count-of-three only as far as A1 carries it; the coproduct earning of addition is S2's.
% EXP-8 placed in S2 (it is the story of HOW the count was earned -- belongs with the earning).
% Gate: turn 510 (Ch2 S1, ~900 words).

An argument that promises to earn every number must begin by saying what may be written
down at all. This chapter speaks grammar, as announced: its claims are rules a reader
can parse, and its objects are texts those rules generate.

Here is the whole ground grammar of the argument:

\begin{center}
\texttt{atom ::= 0 | 1 | 2}
\end{center}

Three marks. Everything in this book that is not proved is generated by this rule or
earned from things that are, and the first discipline of the ground is to take the rule
literally: it is not an abbreviation for "the numbers," not the front end of an infinite
supply. It generates exactly three distinguishable marks, and a mark's whole job is to
be distinguishable --- from the blank page, and from the other two. Nothing else about
an atom is used anywhere in the argument. Not magnitude, not order, not arithmetic:
distinguishability alone.

Why three? Because three is the least ground on which the rest can be earned, and the
ladder that explains why reads directly off the marks. The mark \texttt{0} is a fact:
something stands on the page rather than nothing, and the page can testify to that. The
mark \texttt{1} is a number: with two distinguishable states of affairs --- blank and
marked --- one of them can be counted. The mark \texttt{2} is a natural: two
distinguishable marks give succession its first instance, one-more-than, the seed of
every count that follows. And then the ladder does something different at its next
rung. It does not add a mark.

The numeral 3 is not an atom. It never appears on the right-hand side of the ground
grammar; no rule generates it; the reader will not find it standing on the page as a
mark among marks. It enters the argument in one way only: as the length of the roster.
Write out the language the rule generates --- the whole language, which is finite ---
and you have written a text: \texttt{0}, then \texttt{1}, then \texttt{2}. That text
can be parsed, and one of the things a parse reports is how long the text is. The
report says three. Three is thus a property of a generated text, not a symbol in it: a
fact about the ground, read off the ground, never written into it. The count of three
that Chapter 1 said "governs the bracket's depth" is a read-off length of exactly this
kind --- read, as the next section shows, off a text built from these atoms, not off
this roster itself --- and the reader should hold the book to that, because a 3 that had
been written down by an author, rather than read off a roster, would be contraband under
the first discipline.

This is what the fencepost means, and the surveyor's word is the right one. A fence
with three posts has two spans; the posts are stood on, the spans are walked. In the
ladder of the ground, \texttt{0}, \texttt{1}, \texttt{2} are posts: the grammar stands
each of them on the page. Three is the far end of the survey --- the place the count
points to when the posts run out. It is approached by construction: generate the
roster, parse it, read its length, and you arrive at three every time, with proof. But
it is never landed on: no construction of the ground puts a third-successor mark on the
page, and the argument never behaves as if one were there. Three is a limit in the
plainest sense the grammar affords --- the first numeral the ground can point to but
not touch. Later chapters will make heavy use of exactly this: the procedure that
deepens the bracket runs its count to the fencepost, and stops, because past the
fencepost there is nothing generated to stand on.

The ladder --- fact, number, natural, limit --- should be held the way the argument
holds it: as an interpretation riding on the construction, not as a second
construction. What carries weight is small and checkable: a three-line grammar, its
finite roster, and a length read off a parse. The names on the rungs explain why the
ground was cut here and not elsewhere; they are commentary, and the book will not hang
a theorem on them. If a reader distrusts the ladder entirely, nothing downstream is
owed a revision, so long as the roster and its read-off three are conceded.

What is deliberately not here is nearly everything. There is no four --- the grammar
generates three marks and the argument never pretends otherwise; whatever lies above
the roster is either earned in later chapters or fenced. There is no order among the
atoms yet, and no arithmetic: it is not yet even meaningful to say that \texttt{1} and
\texttt{2} "make" three, because making --- putting texts beside texts and counting
the result --- is an operation, and operations must be earned. That earning is the next
section's single subject, and it comes with a story the ground owes the reader: the
first time this argument tried to obtain its three, it reached for a count that was
quietly circular, and the discipline of this chapter is what caught it.

\section{The earning of addition}
% Node-map: A2 (addition via the coproduct: disjoint alphabets side by side, combined
% roster length read off -- EARNED, never a postulated +), A3 (earn-site discipline;
% universal claim = the audit, PARTIAL-by-nature stated as practice), EXP-8 (the
% non-circularity incident as narrative spine), EXP-6 (the exact-arithmetic audit,
% referenced under A3). Required by turn 511: THE bracket-governing count-3 identified
% with the COPRODUCT roster's read-off length. Gate: turn 511 (Ch2 S2, ~900 words).

The previous section closed owing a story, and the ground pays its debts first.

The argument needed its three --- the count that governs the bracket's depth --- and the
first three it reached for was already lying around. Far downstream of the ground, the
device exhibits a rotation with three positions: a cycle that returns to its start after
three turns, as later chapters will show in detail. Ask how long that cycle takes to
come home and the answer is three; what could be more natural than to let that answer be
the count? But the discipline of the previous section --- every number must name the
text it was read off --- turned the question around. Off which text was this three read?
And the answer pointed downstream, into the running argument, at machinery whose own
credentials rest on the very count being sought. The cycle's period would have certified
the count, and the count would have underwritten the cycle: a circle, quiet, plausible,
and fatal. The gate caught it, and the reach was rejected --- not because the answer
three was wrong, but because it had not been earned. This incident is reported here the
way the fired null was reported in Chapter 1: the discipline is not decoration; it has
drawn blood inside this argument, on this exact number.

The re-earning is the one construction this section performs. Take two alphabets with
nothing in common. The first has two letters; the second has one:

\begin{center}
\texttt{pair ::= a | b} \qquad \texttt{post ::= c}
\end{center}

Now set them side by side --- not merged, not pooled, but placed disjointly, each letter
still wearing the name of the alphabet it came from:

\begin{center}
\texttt{side-by-side ::= pair | post}
\end{center}

Write out the language this generates, and you have written a new text: \texttt{a} from
the left, \texttt{b} from the left, \texttt{c} from the right. Parse it. The parse
reports a length, and the report says three. Nothing was added, in the schoolroom sense,
anywhere in that procedure: no plus sign appears in the grammar, no sum was computed.
A text was built and a length was read, exactly as the atom roster's length was read in
the previous section --- but off a different text, and the difference is the entire
point.

This is addition, earned. "Two beside one is three" is not an appeal to an operation
the argument never justified; it is a parse report on a constructed text. When a plus
sign appears anywhere later in this book, it abbreviates this and only this: set the two
texts side by side, disjointly, and read the combined roster's length. Addition is a
fact about rosters before it is a rule about numerals, and in this argument it is never
anything else.

And here the identity the previous section deferred can be stated where it lives: THE
count of three that governs the bracket's depth is the read-off length of this
side-by-side text --- the two-beside-one roster, built from the atoms, but not the atom
roster itself. Same number; different text; and the bracket hangs on the difference,
because this text's credentials were checked while the tempting one's were not. The
story then closes cleanly: once the count was earned here, the downstream cycle was
proved to match this side-by-side text exactly, letter for letter. The circular
shortcut, once forbidden, was re-derived as a consequence --- the cycle's period of
three became a theorem resting on the earned count, rather than a licence propping it
up. The gate cost the argument nothing but honesty, and bought the ground its floor.

The general discipline, stated once for the whole book: every number in this argument
can be asked, "off which text were you read, or by which proof were you forced?" ---
and has exactly one answer, on file. The ledgers at the argument's close record the
site for each. Honesty requires saying what kind of claim this is: it is an audit, not
a theorem. No master statement certifies all the numbers at once; there is a discipline
applied number by number, and it has been swept --- every quantity in the argument is
an exact ratio of counts, end to end, with nothing rounded and no decimal approximation
anywhere load-bearing. The reader who finds a number in this book without an earn-site
has found a defect, and is asked to report it as one.

The ground now holds atoms and addition. It owes two things more before the instrument
can be introduced: the single sanctioned act of identification --- the one place the
argument is permitted to declare two descriptions the same --- and the discipline of
naming. They are the next sections' subjects, and they complete what may be assumed.

\section{The one identification}
% Node-map: A4 (axiom hygiene GRADED not stated -- apparatus discipline read off the
% instrument's log, METHOD tag kept honest; the needle = the single sanctioned
% identification, described never cited by identifier; the no-choice stance as content:
% every selection the cheapest one, labeled -- plants A5/naming for S4).
% Citations (ledger rows, per apparatus ruling): quotient tradition; no-choice stance.
% Gate: turn 512 (Ch2 S3, ~800 words).

The ground now has marks and it has side-by-side counting, and one subtle power is still
missing. When may the argument say that two different texts are the same thing?

Under the ground's own discipline the question is loaded. A mark's whole job is to be
distinguishable, and everything built so far inherits that default: different texts are
different, full stop. Identification runs against the grain, and it should, because
every act of "these two are the same" erases a distinction, and erased distinctions are
exactly how contraband enters an argument like this one. A smuggled identification can
manufacture a count --- collapse two letters somewhere upstream and every roster
downstream shortens --- so the power to identify is the power to bend every measurement
in the book. It cannot be handed out freely.

The argument's rule is austere: it permits itself exactly one act of identification.
One place --- located, on file, in the apparatus's log, where the book can point to it
--- at which two descriptions that parse differently are declared to name the same
truth. It sits at the one question the argument meets that is genuinely undecidable
from the texts alone: two descriptions, each lawful, each parseable, whose sameness no
finite comparison of letters can settle --- and there, once, the apparatus is permitted
to bring the needle down and declare them one. Everywhere else in the argument, without
exception, sameness is not available by declaration. Two texts that are to be treated
as the same must be proved the same, letter for letter, the way the previous section's
downstream cycle was proved to match the side-by-side roster. Distinguishable until
proved: that is the standing rule, and the one identification is its single, recorded,
priced exception.

Why one, and not zero? Because with no lawful act of identification the book's paired
spellings would be four unrelated monologues. To say that grammar and function-calculus
are two spellings of one content --- to say "same content" at all, anywhere --- the
argument needs one place where sameness-in-truth is lawfully minted. Every use of the
null method downstream, every "subtract the two spellings and read the residue," spends
coin from this one mint. And why one, and not several? Because each further act would
be another place a distinction died unexamined, and the count of such deaths is part of
what the instrument must publish. The published count is one.

That word --- published --- is the section's second subject. The claims of this
paragraph are not theorems of the book, and the reader should notice the difference in
kind. What the argument assumed is a fact about the apparatus, and it is established
the way apparatus facts are established: by reading the instrument's log after the run,
not by arguing beforehand that the log must be clean. The instrument records everything
it was permitted to take for granted, the record is read off it mechanically, and this
book reports what the record shows. It shows two entries. The first is a bookkeeping
principle: two statements proved equivalent may thereafter be written as equal --- the
ledger-keeping form of extensionality, which prices nothing and picks nothing. The
second is the one identification described above. There is no third entry.

And there is a conspicuous absence, which is content in its own right. The log contains
no oracle of choice: no principle anywhere that reaches into an unnamed plurality and
picks a representative without saying which. Nothing in this argument is "some element,
it does not matter which." Every selection made anywhere downstream is the cheapest one
that works, and it is labeled --- picked because it is least, and written down by name.
A pick without a label is contraband of the same species as an unearned number, and the
reader is invited to hunt for one with the same standing invitation the falsifiers
carry. What a label is --- what it costs to name something, and why the cheapest name
is the honest one --- is the last thing the ground owes, and it is the next section's
subject.

The disciplines borrowed here are old, and the ledger at the book's close credits them:
the tradition of minting sameness by quotient, taken with its safety intact, and the
tradition of refusing choice, taken as a vow rather than a hardship. The ground
borrowed both the way a laboratory borrows a balance: instruments, not authorities.

\section{Naming}
% Node-map: A5 (a name names a CLASS, not a text -- identity = distinguishability-class
% identity, lawful under distinguishable-until-proved; naming = decision with cost; the
% cheapest sufficient name is the honest one; cost-minimality displayed AS a fence = the
% anti-continuum fence; the two-box pigeonhole stated SMALL, the tower's wrap explicitly
% deferred to the closure chapter). Citation row: Dirichlet (ledger). Chapter closer.
% Gate: turn 513 (Ch2 S4, ~800 words).

The previous section ended with a promise: every selection in this argument is the
cheapest one that works, and it is labeled. This last section of the ground says what a
label is, and the answer completes what may be assumed.

Begin with what a name may lawfully attach to. Not a single text. Under the standing
rule --- distinguishable until proved --- the argument's grip on its objects goes
exactly as far as its power to tell them apart, and no farther. If two texts cannot be
separated by any test the argument owns, then a name that attached to one and not the
other would itself be doing forbidden work: it would mint a distinction nothing paid
for, the mirror image of the smuggled identification the previous section outlawed. So
a name, to be lawful, attaches to the whole class of texts that the argument's own
tests cannot tell apart. Identity, in this book, is identity of distinguishability
class. That is why naming needs no second needle: a lawful name never declares two
things the same --- sameness-in-truth belongs to the one identification alone --- and
never declares two things different beyond what some test witnesses. It only records
where the tests actually drew their lines.

Recording is the right word, because a name is a decision with a cost. To name a thing
is to run distinguishing tests until it stands separated from everything it must not be
confused with, and then --- this is the discipline --- to stop. The label is the record
of the decisions run, and its cost is their count. The cheapest sufficient name is the
honest one: it contains exactly the decisions the argument made, and no others. A more
expensive name --- one that implies distinctions never actually tested --- claims
unearned separations, contraband of the by-now-familiar species. When this book names
a quantity, the name is a receipt: this many decisions, these ones, and nothing
further. The reader has already seen the practice --- the atoms were named by nothing
but their mutual distinguishability, three decisions, and no talk of magnitude was
allowed to inflate the label.

This discipline, taken seriously, is a boundary of the whole argument, and it is the
ground's one fence:

\begin{fence}
No continuum of names enters this argument. Every name in use is the label of a
finitely decided class, and the supply of such labels is finite at every stage of the
book. Whatever would require infinitely many simultaneous distinctions to name lies
beyond this fence, and the book does not name it --- not as a hardship, but as the
price of every name being a receipt.
\end{fence}

The fence is load-bearing in both directions. Inward, it guarantees that every named
quantity in the book can be audited back to its decisions --- the earn-site discipline
of this chapter's second section, now extended from numbers to names. Outward, it
marks honestly what this argument cannot talk about: anything whose individuation
needs an infinite supply of simultaneous distinctions. The measurements ahead stay
inside the fence by construction, because they are counts, and counts are finite
decision records by birth.

One small fact about the ground's names, stated now because the ending of the book
turns on it. The tests the ground owns return one of two marks --- a test separates,
or it fails to separate --- and so at the ground, the classes that names attach to
sort into exactly two boxes. Two boxes; and any third name must therefore live in a
box already occupied: it shares, whether it likes it or not. That is the pigeonhole,
stated at its smallest, and Dirichlet's ledger row is earned by this paragraph. What
the sharing forces --- how a tower of names, each obliged to room with an earlier one,
is bent by that obligation until the whole structure closes on itself --- is the
closing chapter's subject, and the ground deliberately does not spend it here. The
reader is asked only to hold the small fact: two boxes, and names must share.

The ground is now complete, and its whole inventory can be said in four lines. Three
atoms, whose only property is distinguishability. Addition, as a parse report on texts
set side by side. One identification, on file, priced, and never repeated. And naming,
at cost, every label a receipt. Everything else in this book is measured, proved, or
fenced. The next chapter stands the instrument on this floor: the machine that counts
its own steps, and the discipline that keeps its readout honest.
